\section{Introduction}

Vision-based player tracking offers a scalable alternative to the expensive dedicated hardware traditionally used in sports analytics. 
However, basketball poses severe challenges for these systems \cite{survey}: rapid movements cause frequent mutual occlusions, 
the fast-moving ball suffers from severe motion blur, and indoor environments introduce specular floor reflections. 
Multi-camera setups alleviate occlusions by providing complementary viewpoints, but they require robust cross-view association, 
which in turn depends on consistent player tracking within each individual camera.

In this work, we present an end-to-end pipeline for multi-view player tracking and 3D reconstruction, 
applied to basketball footage from the Sanbapolis facility in Trento, Italy. 
We systematically compare detection strategies: from classical background subtraction (MOG2), to a domain-adapted YOLOv26 
with multi-pass inference—showing that domain adaptation combined with class-agnostic NMS \cite{nms} is essential for reliable detection. 
For tracking, evaluating SORT \cite{sort} against DeepSORT \cite{deepsort} demonstrates that appearance-based re-identification significantly 
reduces identity switches during player crossings. 
Finally, we reconstruct 3D trajectories via SVD-based triangulation across three calibrated cameras, utilizing Kalman smoothing to suppress noise.